% ===== §7 The Practice of the Water-Virtue =====
\section{The Practice of the Water-Virtue: Water-Good Conduct in Intimate Relations}
\label{p20:sec:practice}

\noindent
An ethics that stopped at the criterion would have said what the water-good \emph{is} without saying what it is to \emph{do} it, and for an ethics of intimate relation that would be a failure of nerve, since intimacy is lived in conduct before it is understood in theory. This section carries each water-virtue down into the practice of an actual relation. It is not, however, a manual of advice, and it is built to resist becoming one. Every practice below is derived as an applied corollary of a criterion already established, not offered as a tip; each is marked with its place in the formal scheme; and each is paired with the vice it most easily becomes --- because a practice named only in its good form gives no power to tell itself from its counterfeit, and it is exactly the counterfeits, which wear the look of the virtue, that do the damage. The discipline throughout is threefold: name the good act, name its nearest counterfeit, and give the test that tells them apart. The section closes on the limit of all such telling (\xref{p20:sub:practice-limit}).

\subsection{The practice of non-contention: giving without reclaiming command}
\label{p20:sub:practice-buzheng}

The practice of non-contention is to support the partner's work, growth, and other attachments without converting that support into a claim. One gives --- time, labour, encouragement, the yielding of one's own preference --- and one does not book the giving as a debt the other now owes, nor let it accrue into a standing entitlement to their compliance. This is benefiting that does not reclaim command: the everyday form of generating yet not possessing (\xref{p20:sub:noncon-bushi}). In the formal scheme it is the solenoidal gift, the support that returns to the field of the relation rather than to a ledger held by the giver \parencite{wan2026braid}. Its counterfeit is controlling generosity: the same outward giving, recorded. ``After all I have done for you'' is its signature --- the conversion of accumulated benefaction into leverage, the building of \emph{imperium} out of gifts. The two are identical in their visible conduct and opposite in their structure; what distinguishes them is invisible from outside and decisive from within --- whether the surplus given has been retained as a debt. The test is therefore not how much one gives but what one does with the giving afterward: does the support, once given, stay given, or is it held in reserve as a claim to be called in? A gift kept on the books was never the water-virtue; it was a loan in the costume of a gift.

\subsection{The practice of dwelling-below: flowing to the other's low ground}
\label{p20:sub:practice-chuxia}

The practice of dwelling-below is to move, of one's own accord, toward the partner's least defended ground --- the illness that confers no dignity, the failure with no redemptive arc, the shame, the primal wound, the unbecoming need --- and to be present there without first requiring that the ground be raised. Its mark, established earlier (\xref{p20:sub:low-intimacy}), is that it dwells beside rather than rushing to repair: the companionship that does not need the situation to become other than it is in order to share it \parencite{paperxix}, the epistemic hospitality that makes room for the other's reality including the parts they would hide \parencite{paperxiii}. In the formal scheme it is a flow directed toward the low-potential region, and required to be solenoidal --- generative, returning --- rather than gradient. Its counterfeit is the saviour's descent: a going-down that is really a positioning-above, the rescuer who attends the other's low place in order to occupy the elevated role of the one who rescues, and who therefore needs the other to remain in need. This is dwelling-below as pure gradient --- it descends, but it generates nothing and returns nothing, fixing the other in the dependence from which the saviour draws his standing. The test is the one the formal companion's distinction between power-to and power-over makes sharp: does the dwelling-below leave the other more able to stand, or more dependent on one's standing-by? Care that needs to be needed is not the water-virtue; it is domination that has learned to look like tenderness.

\subsection{The practice of non-fullness: leaving room, not seeing-through}
\label{p20:sub:practice-buying}

The practice of non-fullness is to leave the partner room to be more, and other, than one has known --- to hold one's knowledge of them open, to meet today someone who may have changed since yesterday, to decline to settle them into a finished character. It is the refusal to complete one's reading of the other (\xref{p20:sub:nonfull-intimacy}), grounded in the principle that the other is never exhausted \parencite{paperxix}. In the formal scheme it is the maintenance of the relation away from the solidified fixed point, the refusal to let accumulated history harden into a fixed channel \parencite{wan2026braid}. Its counterfeit is the cognitive domination that travels under ``I know you best'': the use of one's knowledge of the partner not to meet them but to fix them --- to pre-empt their words, to predict their reactions back at them, to deny them the standing to become someone the prediction did not foresee. This is solidification wearing the mask of intimacy; it offers completeness of knowledge as closeness while using that completeness to freeze the other in place. The test is whether one's knowing leaves the other free to change: does ``I know you'' open room for who they are becoming, or nail them to who they have been? Knowledge that forecloses the other's revision is not intimacy's depth; it is the relation setting into the shape it will keep until it breaks.

\subsection{The practice of not-ruling: not shaping the other's becoming}
\label{p20:sub:practice-buzai}

The practice of not-ruling is to let the partner grow after their own kind rather than bending them toward the shape one prefers --- and its difficulty is that it can be hard to distinguish, from inside, from a genuine wish for the other's good. The distinction it requires is between helping the other become a better version of themselves and helping them become a version more to one's own liking. The first serves a growth whose direction is the other's own; the second annexes the other's becoming to one's preference and calls the annexation love. In the formal scheme this is the respect for what the companion paper locates as the flesh --- the non-substitutable particularity of \emph{this} person, which no specification of a preferred type can capture \parencite{wan2026braid} --- against the substitution of the skeleton, the conformity to a type, for the token actually loved. Its counterfeit is remaking love: the love whose real object is an ideal version of the partner, toward which the actual partner is treated as a draft to be corrected. The lover who remakes is contending over the other's becoming in the precise sense developed earlier (\xref{p20:sub:noncon-intimacy}); he possesses by shaping. The test cuts to the root of what is loved: do you love him, or the version of him you are cultivating? A love directed at who the other ought to become has not yet met who the other is --- and it is only the second that can be loved as a person rather than as a project.

\subsection{The unifying criterion: every practice holds the rigid lower bound}
\label{p20:sub:practice-bound}

The four practices share a single direction of failure, and naming it is what binds them into one ethic rather than four counsels. Each can collapse, on the side of softness, into self-dissolution: non-contention into a boundary-less self-emptying that gives everything and reserves no self; dwelling-below into a masochistic submission that takes the lowest place because it believes it deserves no other; non-fullness into the permitting of harm, the room-leaving that becomes a refusal to ever hold the other to account; not-ruling into an unprincipled indulgence that calls the abdication of all shaping a respect for the other's freedom. Each of these is the gradient collapse of its virtue --- the yielding that has fallen below the rigid lower bound (\xref{p20:sub:dialectic-bound}). The unifying criterion of the practice is therefore the one the dialectical core established: water-good conduct is not unconditional yielding but yielding that holds the bound --- that circulates and returns, that keeps the boundary refusing to be drained as fuel. In the register of practice this is what water's other face means. The same water that takes the vessel's shape wears through stone and breaches the dyke; in conduct, this is the requirement that when a relation begins to draw one off as mere fuel --- when the flow has become all one way, the surplus all retained on the other side, one's own measure collapsing --- the water-virtue calls not for further yielding but for the breach: for the force that wears through, or the flood that overruns the bank, or simply the departure of the water from a vessel that has become a trap. This is the decisive answer, at the level of lived practice, to the feminist challenge the survey raised \parencite{wan2024grb}: the water-virtue never asked anyone to remain in a relation that dissolves them. To stay and be drained is not the practice of the water-good; it is the practice fallen below its bound, and the bound itself requires that one wear through or break out. The softness has an edge, and the edge is for exactly the case where softness would otherwise become one's ruin.

\subsection{The limit of practice: the water-good cannot be made a procedure}
\label{p20:sub:practice-limit}

A final honesty is owed, and it guards this section against the very thing it might be mistaken for. Nothing above is a checklist, and the tests offered are not algorithms that decide the cases for us. The formal companion gives the reason in exact form: it proves that a relational structure powerful enough to decide every case --- to compute any distinction whatever --- thereby loses the capacity to discriminate good circulation from bad, so that the very omnipotence one would want in a complete procedure is incompatible with the discriminating power a normative criterion requires \parencite{wan2026braid}. Read into the practice, this means there is no complete mechanical procedure for the water-good; a rule-set fine enough to settle every case in advance would, for that reason, have lost the ability to tell the water-virtue from its counterfeit. Each occasion must therefore be judged afresh, and judged in the flesh --- in the non-substitutable particularity of \emph{this} relation, which no general rule reaches, and which the formal scheme itself locates as the remainder no criterion can compute (\xref{p20:sub:formal-flesh}, \xref{p20:sub:formal-incompleteness}). The practice of the water-virtue thus avows its own incompleteness, and the avowal is not a weakness but a consequence: of a piece with the epistemic humility that has marked this work from its founding \parencite{wan2024grb}, it is the practical form of the recognition that the good of an intimate relation is not the kind of thing a procedure can contain. The tests above are not the substitutes for judgment but its instruments; they sharpen the question each occasion poses, and they leave the answering of it where it belongs --- in the hands of the two who are in the relation, facing this case, which has never occurred before and will not recur.
